Algorithmic redistricting replaces human mapmakers with deterministic computation,
but algorithms can still be gamed if the adversary has freedom to choose input
data, parameters, geographic definitions, which results to present, or which
algorithm version to deploy.
We systematically enumerate five gaming vectors for the \bisect{} redistricting
platform and measure the maximum achievable partisan shift under each vector
individually and in combination.

The five vectors are: (V1)~parameter manipulation, (V2)~input data manipulation,
(V3)~geographic definition gaming, (V4)~selective result presentation, and
(V5)~algorithm version selection.
Vectors V1--V3 affect the generated plan itself; V4 and V5 affect only how
results are reported, not the underlying plan geometry.

Empirical analysis across all 50 states reveals that the combined maximum
partisan effect of all plan-affecting vectors (V1--V3) is at most 0.5
Democratic seats nationally under additive worst-case assumptions and
at most 0.3 seats in practice.
This is comparable to the seed-noise floor established in B.7 (CV $<$ 2\%)
and far smaller than the 18-seat gap between algorithm structures documented
in B.0.

The \dia{} pre-registration protocol eliminates all five vectors for
certified plans: parameters are fixed in a SHA-256 formula before
Census data are released, input data are verified against Census Bureau
hashes, geographic definitions are mandated by statute, result presentation
follows a mandatory disclosure standard, and the algorithm version is
pinned in the manifest.
We conclude that when the \dia{} protocol is followed, no gaming vector
can produce an undetectable partisan outcome larger than the algorithm's
inherent sampling noise.
